\documentclass[12pt,a4paper]{article}
\usepackage[utf8]{inputenc}
\usepackage[T1]{fontenc}
\usepackage[brazil]{babel}
\usepackage{amsmath,amssymb}
\usepackage{geometry}
\geometry{margin=2.5cm}
\title{Material de Estudo: Apresenta\c{c}\~ao de Trabalhos II --- Comunica\c{c}\~ao e Defesa}
\author{Princ\'ipio de Modelagem Matem\'atica}
\date{Unidade 33}
\begin{document}
\maketitle

\section*{Objetivo}
Desenvolver habilidades de comunica\c{c}\~ao oral cient\'ifica, incluindo como responder perguntas da audi\^encia.

\section*{Roteiro de Apresenta\c{c}\~ao (10 min)}

\begin{enumerate}
  \item \textbf{0--1 min:} Apresente-se, o t\'opico e por que ele \'e importante.
  \item \textbf{1--2 min:} Contexto e pergunta de pesquisa.
  \item \textbf{2--4 min:} O modelo: equa\c{c}\~oes, hip\'oteses, par\^ametros.
  \item \textbf{4--7 min:} Resultados principais (2--3 figuras chave).
  \item \textbf{7--9 min:} Discuss\~ao, limita\c{c}\~oes, compara\c{c}\~ao.
  \item \textbf{9--10 min:} Conclus\~ao e perspectivas.
  \item \textbf{+3 min:} Perguntas.
\end{enumerate}

\section*{Como Responder Perguntas}

\begin{itemize}
  \item \textbf{Boa pergunta:} repita a pergunta (ajuda todos a ouvir), depois responda.
  \item \textbf{N\~ao sabe a resposta:} ``Essa \'e uma excelente quest\~ao. N\~ao investigamos isso, mas uma possibilidade \'e\ldots'' \'e melhor que inventar.
  \item \textbf{Pergunta fora do escopo:} ``Isso est\'a al\'em do que modelamos, mas \'e uma extens\~ao interessante.''
  \item \textbf{Cr\'itica:} ouça com aten\c{c}\~ao, reconhe\c{c}a o ponto, explique suas escolhas.
  \item \textbf{Perguntas simuladas para praticar:}
  \begin{itemize}
    \item ``Por que voc\^e escolheu esse m\'etodo num\'erico e n\~ao outro?''
    \item ``Como voc\^e validou seu modelo?''
    \item ``O que mudaria se a hip\'otese X fosse relaxada?''
    \item ``Quais s\~ao as limita\c{c}\~oes mais cr\'iticas?''
  \end{itemize}
\end{itemize}

\section*{Comunica\c{c}\~ao Visual}

\subsection*{Figuras eficazes}
\begin{itemize}
  \item Eixos com r\'otulos e unidades sempre.
  \item T\'itulo informativo: ``MSD vs.\ tempo para LJ a $T^*=1{,}5$'' n\~ao s\'o ``MSD''.
  \item Legenda dentro da figura quando poss\'ivel.
  \item Escala coerente: n\~ao distorcer.
  \item Cores acess\'iveis (evite vermelho-verde para daltonismo).
\end{itemize}

\subsection*{Equa\c{c}\~oes}
\begin{itemize}
  \item Mostre apenas as equa\c{c}\~oes centrais do modelo.
  \item Explique cada s\'imbolo na hora que aparecer.
  \item Nunca ponha uma equa\c{c}\~ao no slide sem discuti-la.
\end{itemize}

\section*{Exerc\'icios}

\begin{enumerate}
  \item \textbf{(Ensaio)} Prepare uma apresenta\c{c}\~ao de 2 minutos sobre seu tema. Grave em v\'ideo, assista e identifique 3 pontos a melhorar.
  
  \item \textbf{(Perguntas)} Para cada pergunta simulada acima, escreva uma resposta de 3--5 frases adequada para o seu modelo.
  
  \item \textbf{(Figura)} Pegue um gr\'afico de resultado do seu trabalho. Verifique: (a) h\'a r\'otulos em todos os eixos? (b) a unidade est\'a indicada? (c) o t\'itulo \'e informativo? (d) a legenda \'e clara? Corrija o que faltar.
  
  \item \textbf{(Tempo)} Cronometre cada se\c{c}\~ao da sua apresenta\c{c}\~ao. Quanto tempo voc\^e gasta em cada uma? Est\'a coerente com o roteiro proposto? O que \'e preciso cortar ou expandir?
  
  \item \textbf{(Limita\c{c}\~oes)} Escreva um par\'agrafo de ``discuss\~ao de limita\c{c}\~oes'' do seu modelo. Seja espec\'ifico: n\~ao s\'o ``o modelo \'e simplificado'', mas ``assumimos X, o que ignora Y, e isso pode causar Z''.
  
  \item \textbf{(Contexto)} Pesquise 2--3 trabalhos relacionados ao seu tema. Como seu modelo se diferencia? Quais avanços voc\^e apresenta?
  
  \item \textbf{(Perspectivas)} Escreva 3 extens\~oes concretas do seu trabalho: o que faria se tivesse mais tempo? Quais dados coletaria? Qual hip\'otese relaxaria primeiro?
\end{enumerate}

\section*{Checklist para a Apresenta\c{c}\~ao}
\begin{itemize}
  \item[$\square$] Slides revisados (ortografia, figuras, equa\c{c}\~oes).
  \item[$\square$] Apresenta\c{c}\~ao ensaiada ao menos 2 vezes.
  \item[$\square$] Tempo cronometrado e dentro do limite.
  \item[$\square$] Respostas para perguntas prov\'aveis preparadas.
  \item[$\square$] Arquivo backup dispon\'ivel (nuvem ou pen drive).
\end{itemize}

\section*{Resumo R\'apido}
\begin{itemize}
  \item Roteiro: 7 etapas em 10 min; perguntas nos 3 min finais.
  \item Responder perguntas: repita, seja honesto, reconhe\c{c}a limita\c{c}\~oes.
  \item Figuras: r\'otulos, unidades, t\'itulo informativo.
\end{itemize}

\section*{Refer\^encias}
\begin{itemize}
  \item Notas de aula da disciplina.
\end{itemize}

\end{document}
